% =========================================================================
%  LABORATORY RECORD INDEX
% =========================================================================
\thispagestyle{plain}
\pdfbookmark[0]{Index}{index}

\begin{center}
    {\LARGE\bfseries INDEX}
\end{center}

\vspace{0.5cm}

\renewcommand{\arraystretch}{1.3}
\noindent\begin{tabularx}{\linewidth}{|c|c|>{\raggedright\arraybackslash}X|c|c|c|}
\hline
\textbf{Exp.} & \textbf{Date} & \textbf{Title of Experiment} & \textbf{Page} & \textbf{Marks} & \textbf{Faculty Sign} \\ \hline
1  & -- & Introduction to UNIX/Linux and Basic Commands & \pageref{exp:unix} & & \\ \hline
2  & -- & UNIX System Calls & \pageref{exp:syscalls} & & \\ \hline
3  & -- & Simulation of UNIX Commands (ls, grep) & \pageref{exp:simulate} & & \\ \hline
4  & -- & Shell Programming & \pageref{exp:shell} & & \\ \hline
5  & -- & CPU Scheduling Algorithms & \pageref{exp:scheduling} & & \\ \hline
6  & -- & Producer--Consumer Problem using Semaphores & \pageref{exp:producerconsumer} & & \\ \hline
7  & -- & IPC using Shared Memory & -- & & \\ \hline
8  & -- & Banker's Algorithm for Deadlock Avoidance & \pageref{exp:bankers} & & \\ \hline
9  & -- & Deadlock Detection Algorithm & -- & & \\ \hline
10 & -- & Threading and Synchronization & -- & & \\ \hline
11 & -- & Memory Allocation: First, Worst and Best Fit & -- & & \\ \hline
12 & -- & Paging Technique of Memory Management & -- & & \\ \hline
13 & -- & Page Replacement: FIFO, LRU and LFU & -- & & \\ \hline
14 & -- & File Organization: Single and Two-Level Directory & -- & & \\ \hline
15 & -- & File Allocation: Sequential, Linked and Indexed & -- & & \\ \hline
\end{tabularx}
